\section{Notation and physical units}\label{sec:notation}
Table~\ref{tab:notation} collects the mathematical quantities used in the model,
algorithms and acceptance calculations. A single scalar objective unit is assumed
within one optimization instance. This permits comparison of objective values
between equivalent formulations of that instance, but does not permit adding
objectives from different physical models or independently normalized examples.
The sparse matrix entries inherit row-unit divided by decision-unit dimensions.
Their numerical magnitude may change under unit conversion even when the physical
problem and its optimizer remain identical.

\begin{table}[htbp]\centering\small
\caption{Notation and dimensions for the original-space soft quadratic model. Row and decision units belong to the physical metadata of each synthetic case rather than to a universal unit convention.}\label{tab:notation}
\begin{tabularx}{\textwidth}{@{}lXX@{}}\toprule
Symbol & Meaning & Units\\\midrule
$x_j$ & Original decision component & Declared decision unit\\
$A_i x,\ell_i,u_i$ & Row value and interval endpoints & Declared row unit\\
$v_{L,i},v_{U,i},s_i$ & Lower/upper departure and corrective slack & Row unit\\
$f,F,g_i$ & Base, total and interval penalty objective & Objective unit\\
$a_{L,i},a_{U,i}$ & Linear departure prices & Objective unit per row unit\\
$b_{L,i},b_{U,i}$ & Quadratic departure prices & Objective unit per squared row unit\\
$c_{L,i},c_{U,i}$ & Directional peak scales & Peak unit per row unit\\
$y_i$ & Original row multiplier & Objective unit per row unit\\
$H_{ii}$ & Positive proximal metric entry & Objective unit per squared row unit\\
$r_{\rm stat}$ & Stationarity residual & Coordinatewise objective gradient unit\\
$T,\Delta t$ & Measured duration and physical time step & Seconds, or explicit model hours\\
$R_i$ & Paired relative timing reduction & Dimensionless; displayed as percent\\\bottomrule
\end{tabularx}
\end{table}

A row transformation $z_i'=d_i z_i$ with $d_i>0$ transforms both endpoints by the
same factor. To preserve a separable objective, the corresponding prices transform
as $a_i'=a_i/d_i$ and $b_i'=b_i/d_i^2$. A directional peak scale transforms as
$c_i'=c_i/d_i$. The row dual transforms reciprocally, $y_i'=y_i/d_i$, so its
contribution to stationarity remains invariant. These rules distinguish a genuine
unit change from a policy change. Altering only a directional scale changes the
penalty's relative preference and can legitimately move the physical optimizer.

The control examples use position in metres, velocity in metres per second and
acceleration in metres per squared second. A metre-to-centimetre transformation
must also transform position-sensitive objective and row coefficients. A seconds-
to-milliseconds transformation changes discrete dynamic coefficients and the
numerical representation of velocity and acceleration; changing labels alone is
not a test of covariance. The hourly thermal grouped examples use a different
time convention from the quarter-hour published building example. Their heater
energy and comfort scales are interpreted within their own model, not transferred
by matching variable names.

Timing quantities are wall-clock durations with explicit measurement boundaries.
They are not physical prediction times. A thirty-two-frame forecast operation
can have one measured duration even though its model horizon spans hours or
seconds. Likewise, a plotted input index is a population ordering, not a measured
time axis. Keeping these meanings distinct prevents a computational scatter plot
from being read as a physical control trajectory or an operational latency guarantee.

The unit-aware reader distinguishes a zero value from an absent physical
quantity. If the selected row unit contains no hard row, no maximum hard-row
error in that unit exists. Reporting zero would incorrectly suggest a measured
feasibility check. The same convention applies to unsupported method outputs:
their objective and elapsed duration remain unavailable rather than receiving
a numerical value that could enter a later sum or comparison.
